\chapter{Working with Recursion on Lists}

\tikzset{
  lcell/.style={draw,minimum width=7mm,minimum height=6mm,inner sep=1pt},
  lptr/.style={-{Stealth}},
  rarr/.style={-{Stealth},thick},
  rcall/.style={draw,rounded corners,fill=blue!6,font=\ttfamily\small,minimum height=6mm},
  rval/.style={draw,rounded corners,fill=green!8,font=\ttfamily\small,minimum height=6mm,minimum width=8mm}
}

Having seen the list template, we now practise \textbf{structural recursion}:
recursion whose shape follows exactly the shape of the data definition.

\begin{definition}[Pure structural recursion]
A function uses \textbf{pure structural recursion} if every recursive call is
on one of the self-referential parts of the data (for a list, \code{(rest lst)})
and all other parameters are passed unchanged.
\end{definition}

\section{Functions producing Booleans}

\begin{lstlisting}[style=racket]
;; (contains? item lst) produces true if item is in lst.
;; Examples:
(check-expect (contains? 'b '(a b c)) true)
(check-expect (contains? 'z empty) false)

;; contains?: Any (listof Any) -> Bool
(define (contains? item lst)
  (cond [(empty? lst) false]
        [else (or (equal? item (first lst))
                  (contains? item (rest lst)))]))
\end{lstlisting}
Because \code{or} short-circuits, the recursion stops as soon as the item is
found. (Built-in: \code{member?}.)

\section{Functions producing lists}

\subsection{Transforming each element (a ``map'' pattern)}
\begin{lstlisting}[style=racket]
;; (double-all lon) produces a list with each number in lon doubled.
;; Examples:
(check-expect (double-all empty) empty)
(check-expect (double-all '(1 2 3)) '(2 4 6))

;; double-all: (listof Num) -> (listof Num)
(define (double-all lon)
  (cond [(empty? lon) empty]
        [else (cons (* 2 (first lon)) (double-all (rest lon)))]))
\end{lstlisting}

\subsection{Selecting some elements (a ``filter'' pattern)}
\begin{lstlisting}[style=racket]
;; (keep-positive lon) produces the positive numbers of lon, in order.
;; Examples:
(check-expect (keep-positive '(3 -1 0 5)) '(3 5))

;; keep-positive: (listof Num) -> (listof Num)
(define (keep-positive lon)
  (cond [(empty? lon) empty]
        [(positive? (first lon)) (cons (first lon) (keep-positive (rest lon)))]
        [else (keep-positive (rest lon))]))
\end{lstlisting}

Element by element, the two patterns look like this: \code{double-all}
sends every element through the same operation, while \code{keep-positive}
lets some elements through and drops the rest.

\begin{center}
\begin{tikzpicture}[x=8mm]
  \foreach \v [count=\i] in {1,2,3} {
    \node[lcell] (t\i) at (\i,0) {\v};
    \pgfmathtruncatemacro{\d}{2*\v}
    \node[lcell,fill=green!8] (b\i) at (\i,-1.6) {\d};
    \draw[rarr] (t\i) -- node[right,font=\scriptsize,inner sep=1pt]{$\times 2$} (b\i);
  }
  \node[font=\scriptsize\ttfamily] at (2,0.65) {double-all};
  \foreach \v/\k [count=\i] in {3/1,-1/0,0/0,5/1} {
    \node[lcell] (u\i) at (\i+5,0) {\v};
  }
  \node[font=\scriptsize\ttfamily] at (7.5,0.65) {keep-positive};
  \node[lcell,fill=green!8] (k1) at (6.5,-1.6) {3};
  \node[lcell,fill=green!8] (k2) at (8.5,-1.6) {5};
  \draw[rarr] (u1) -- (k1);
  \draw[rarr] (u4) -- (k2);
  \foreach \i in {2,3} \draw[red!60!black,thick] ($(u\i.center)+(-0.25,-0.5)$) -- ++(0.5,-0.4) ($(u\i.center)+(0.25,-0.5)$) -- ++(-0.5,-0.4);
\end{tikzpicture}
\end{center}

These patterns recur so often that Chapter~\ref{ch:abstraction} abstracts them
into the built-in higher-order functions \code{map} and \code{filter}.

\subsection{Appending}
\begin{lstlisting}[style=racket]
;; (my-append l1 l2) produces the elements of l1 followed by those of l2.
;; my-append: (listof Any) (listof Any) -> (listof Any)
(check-expect (my-append '(1 2) '(3 4)) '(1 2 3 4))
(define (my-append l1 l2)
  (cond [(empty? l1) l2]
        [else (cons (first l1) (my-append (rest l1) l2))]))
\end{lstlisting}
Only \code{l1} is recursed on; \code{l2} goes along unchanged. This is still
pure structural recursion on \code{l1}.

\section{Non-empty lists}
Some functions make no sense on \code{empty}, such as the maximum:
\begin{lstlisting}[style=racket]
;; (list-max lon) produces the largest number in lon.
;; list-max: (ne-listof Num) -> Num
(check-expect (list-max '(3 9 2)) 9)
(check-expect (list-max '(4)) 4)
(define (list-max lon)
  (cond [(empty? (rest lon)) (first lon)]
        [else (max (first lon) (list-max (rest lon)))]))
\end{lstlisting}
The template for \code{(ne-listof X)} has base case \code{(empty? (rest lst))}.

\begin{aside}
A version that computes \code{(list-max (rest lon))} \emph{twice} (once in a
comparison and once as the answer) takes exponential time: each call makes two
recursive calls. Avoid repeating recursive calls; use \code{max}, a helper, or
\code{local} (Chapter~\ref{ch:local}).
\end{aside}

\section{Insertion sort}

Sorting is a classic application of structural recursion. If we can sort the
rest of the list, we only need to \emph{insert} the first element into the
correct place.

\begin{lstlisting}[style=racket]
;; (sort lon) produces lon sorted in non-decreasing order.
;; Examples:
(check-expect (sort empty) empty)
(check-expect (sort '(3 1 2)) '(1 2 3))

;; sort: (listof Num) -> (listof Num)
(define (sort lon)
  (cond [(empty? lon) empty]
        [else (insert (first lon) (sort (rest lon)))]))

;; (insert n slon) inserts n into the sorted list slon so that the
;;   result is sorted.
;; Examples:
(check-expect (insert 4 '(1 3 5)) '(1 3 4 5))
(check-expect (insert 0 empty) '(0))

;; insert: Num (listof Num) -> (listof Num)
;;   requires: slon is sorted in non-decreasing order
(define (insert n slon)
  (cond [(empty? slon) (cons n empty)]
        [(<= n (first slon)) (cons n slon)]
        [else (cons (first slon) (insert n (rest slon)))]))
\end{lstlisting}

\code{insert} only rebuilds the cells in front of the insertion point; the
rest of \code{slon} is shared with the result, not copied:

\begin{center}
\begin{tikzpicture}
  \node[font=\scriptsize\ttfamily,anchor=east] at (-0.6,0) {slon};
  \node[lcell] (a1) at (0,0) {1}; \node[lcell] (a2) at (0.7,0) {};
  \node[lcell] (b1) at (2.1,0) {3}; \node[lcell] (b2) at (2.8,0) {};
  \node[lcell] (c1) at (4.2,0) {5}; \node[lcell] (c2) at (4.9,0) {};
  \node[font=\ttfamily\small] (e) at (6.4,0) {empty};
  \draw[lptr] (a2.center) -- (a2.center -| b1.west);
  \draw[lptr] (b2.center) -- (c1.west);
  \draw[lptr] (c2.center) -- (e.west);
  \begin{scope}[yshift=-1.5cm]
  \node[font=\scriptsize\ttfamily,anchor=east] at (-0.6,0) {(insert 4 slon)};
  \node[lcell,fill=green!8] (x1) at (0,0) {1}; \node[lcell,fill=green!8] (x2) at (0.7,0) {};
  \node[lcell,fill=green!8] (y1) at (2.1,0) {3}; \node[lcell,fill=green!8] (y2) at (2.8,0) {};
  \node[lcell,fill=green!8] (z1) at (4.2,0) {4}; \node[lcell,fill=green!8] (z2) at (4.9,0) {};
  \end{scope}
  \draw[lptr] (x2.center) -- (y1.west);
  \draw[lptr] (y2.center) -- (z1.west);
  \draw[lptr,thick,blue!60!black] (z2.center) -- ++(0.9,0) |- ($(c1.south)+(0,-0.25)$) -- (c1.south);
  \node[font=\scriptsize,text=blue!60!black,anchor=west] at (5.9,-0.8) {shared tail};
\end{tikzpicture}
\end{center}

\paragraph{Condensed trace.}
\begin{lstlisting}[style=racket]
(sort (list 3 1 2))
=> (insert 3 (sort (list 1 2)))
=> (insert 3 (insert 1 (sort (list 2))))
=> (insert 3 (insert 1 (insert 2 (sort empty))))
=> (insert 3 (insert 1 (insert 2 empty)))
=> (insert 3 (insert 1 (list 2)))
=> (insert 3 (list 1 2))
=> (cons 1 (insert 3 (list 2)))
=> (cons 1 (cons 2 (insert 3 empty)))
=> (list 1 2 3)
\end{lstlisting}

The recursion first unwinds all the way to \code{empty}; the sorted list is
then built on the way back up by one \code{insert} per level:

\begin{center}
\begin{tikzpicture}
  \node[rcall] (c1) at (0,0) {(sort (list 3 1 2))};
  \node[rcall] (c2) at (0,-1.1) {(sort (list 1 2))};
  \node[rcall] (c3) at (0,-2.2) {(sort (list 2))};
  \node[rcall] (c4) at (0,-3.3) {(sort empty)};
  \node[rval] (r4) at (6,-3.3) {empty};
  \node[rval] (r3) at (6,-2.2) {(list 2)};
  \node[rval] (r2) at (6,-1.1) {(list 1 2)};
  \node[rval] (r1) at (6,0) {(list 1 2 3)};
  \foreach \a/\b in {c1/c2,c2/c3,c3/c4} \draw[rarr] (\a) -- (\b);
  \draw[rarr] (c4) -- (r4);
  \draw[rarr,green!40!black] (r4) -- node[right,font=\scriptsize\ttfamily]{insert 2} (r3);
  \draw[rarr,green!40!black] (r3) -- node[right,font=\scriptsize\ttfamily]{insert 1} (r2);
  \draw[rarr,green!40!black] (r2) -- node[right,font=\scriptsize\ttfamily]{insert 3} (r1);
  \foreach \a/\b in {c1/r1,c2/r2,c3/r3} \draw[densely dotted] (\a) -- (\b);
\end{tikzpicture}
\end{center}

Insertion sort takes time proportional to $n^2$ in the worst case (a
reverse-sorted list), since each of $n$ insertions may traverse the whole
sorted list. (The name \code{sort} clashes with a built-in in later levels;
call it \code{insertion-sort} there.)

\section{Lists of lists}
A list element may itself be a list. With nested data, the template calls a
helper (or another template) on \code{(first lst)}:
\begin{lstlisting}[style=racket]
;; (sum-all lolon) produces the sum of all numbers in all lists of lolon.
;; sum-all: (listof (listof Num)) -> Num
(check-expect (sum-all '((1 2) () (3))) 6)
(define (sum-all lolon)
  (cond [(empty? lolon) 0]
        [else (+ (sum (first lolon)) (sum-all (rest lolon)))]))
\end{lstlisting}

\section{Wrapper functions}
A \textbf{wrapper} is a short non-recursive function that adapts its input and
calls a recursive helper, e.g.\ \code{count-vowels} (Chapter~\ref{ch:lists})
converts a string to a list of characters and calls the list function. The
wrapper has the nice interface; the helper does the recursion.

\section{Testing recursive functions}
Test at least: the empty list; a one-element list; a longer list; lists where
the interesting element is first, last, absent, or repeated; already sorted
and reverse-sorted input for sorting functions.

\section{Summary}
\begin{itemize}
  \item Structural recursion follows the data definition; one recursive call on \code{(rest lst)}.
  \item Common patterns: count/sum (combine), contains (search with \code{or}), map (transform each), filter (keep some).
  \item Non-empty lists have a base case at a one-element list.
  \item Insertion sort: sort the rest, then insert the first; $\bigO(n^2)$ worst case.
  \item Never repeat the same recursive call in one body.
\end{itemize}
